\documentclass[11pt]{paper}

\usepackage[T1]{fontenc}

\usepackage{graphicx}
\usepackage{amssymb}
\usepackage{amstext}
\usepackage{amsmath}
\usepackage{a4wide,color}
\usepackage[utf8]{inputenc}
\usepackage[frenchb]{babel}
\usepackage{xspace}
\usepackage{anysize}
\usepackage{tabularx}
\usepackage{multirow}
\usepackage{fancybox}
\usepackage{fancyhdr}
\usepackage{bbding}

\usepackage{color}
\usepackage{float}

\usepackage[toc,page]{appendix} 
\usepackage{lscape}

\usepackage{placeins}

\usepackage{listings}
\lstdefinestyle{customc}{
  belowcaptionskip=1\baselineskip,
  breaklines=true,
  frame=L,
  xleftmargin=\parindent,
  language=C,
  showstringspaces=false,
  basicstyle=\tiny\ttfamily,
  keywordstyle=\bfseries\color{green!40!black},
  commentstyle=\itshape\color{purple!40!black},
  identifierstyle=\color{blue},
  stringstyle=\color{orange},
  tabsize=2
}

\usepackage{todonotes}

\newcommand{\version}{1.0.1{ }}

\usepackage[framemethod=TikZ]{mdframed}

\lstset{% general command to set parameter(s)
	basicstyle=\footnotesize, % print whole listing small
	keywordstyle=\color{magenta}\bfseries, % underlined bold black keywords
	identifierstyle=, % nothing happens
	commentstyle=\color{black}, % white comments
	stringstyle=\ttfamily, % typewriter type for strings
	showstringspaces=false % no special string spaces
} 

\lstdefinelanguage{aald}
	{morekeywords={system, implementation},
	sensitive=false,
	morecomment=[l]{//},
	morecomment=[s]{/*}{*/},
	morestring=[b]",
}

\usepackage[colorlinks=true]{hyperref}

\usepackage[normalem]{ulem}
\usepackage{color}

\definecolor{Fond}{gray}{0.7}

\renewcommand{\floatpagefraction}{.99}
\renewcommand{\textfraction}{.01}
\newcommand{\modif}[1]{\textcolor{red}{\uline{#1}}}


\pagestyle{fancy}
\fancyhf{}
\fancyhead[RE,RO]{\thepage}
\fancyhead[LE]{}
\fancyhead[LO]{Temps réel }
\fancyfoot[LO]{Centrale Nantes -- \today}
\fancypagestyle{plain}{%
  \fancyhf{} % get rid of headers
  \renewcommand{\headrulewidth}{0pt} % and the line
}

\newenvironment{maliste}%
{ \begin{list}%
	{\ArrowBoldRightStrobe}%
	{\setlength{\labelwidth}{30pt}%
	 \setlength{\leftmargin}{35pt}%
	 \setlength{\itemsep}{\parsep}}}%
{ \end{list} }

 \setlength{\parskip}{1.2ex plus .9ex minus .4ex}

\newcommand{\ocv}{{\tt OpenCV}\xspace} 
\newcommand{\alsa}{{\tt ALSA}\xspace} 
\newcommand{\aadl}{{\tt AADL}\xspace} 
\newcommand{\xeno}{{\tt Xenomai}\xspace} 
\newcommand{\iic}{I$^2$C\xspace} 


\title{{\Large Systèmes interconnectés -- SINT }\\
\vspace{-0.4em}
{\large Sujet de TP 3 : SPI}\\
{\small Centrale Nantes}\\
\small P.-E. Hladik, \small{\texttt{pehladik@ec-nantes.fr}}\\
---\\
{\scriptsize  Version \version (\today)}\\
}


\begin{document}

\maketitle

%\begin{table}[htdp]
%\caption{Suivi des modifications}
%\begin{center}
%\begin{tabular}{|c|c|c|p{8cm}|}
%\hline
%Date & Version & Auteur & Modifications\\
%\hline
%JJ/MM/12 & 0.3.X & P.-E. Hladik&  \\
%\hline
%\end{tabular}
%\end{center}
%\label{default}
%\end{table}%

%%%%%%%%%%%%%%%%%%%%%%%%%
  \section{Objectifs pédagogiques}
%%%%%%%%%%%%%%%%%%%%%%%%%


\begin{itemize}
	\item écriture d'un driver logiciel pour le SPI
	\item écriture d'un driver de communication MCP2515
\end{itemize}


%%%%%%%%%%%%%%%%%%%%%%%%%
  \section{Ce qu'il y a à rendre}
%%%%%%%%%%%%%%%%%%%%%%%%%

Vous déposez sur hippocampus un document pdf avec les parties pertinente de votre code commenté.

%%%%%%%%%%%%%%%%%%%%%%%%%
  \section{Ecrire un drive logiciel}
%%%%%%%%%%%%%%%%%%%%%%%%%

Ce TP s’effectue sur la carte Feather M0 WiFi. Vous pouvez piocher dans le code des deux activités précédentes.

Ce TP  consiste à programmer une liaison {\bf SPI logicielle} pour communiquer avec le MCP2515. Vous ne devez pas utiliser le périphérique SPI mais uniquement communiquer via les ports d'entrée-sortie (voir le cours). Le datasheet du MCP2525 est disponible sur le serveur pédagogique (a priori ça ne vous servira pas, tout a déjà été dit en cours). 

Vous devez produire quatre fonctions :
\begin{itemize}
	\item {\bf Initialisation :} configuration des ports d'entrée-sortie.\\
	{\tt void initialiserSPILogiciel (void) ;}
	
	\item {\bf  Échange de données SPI :} Cette fonction est appelée par les deux fonctions suivantes. Elle sérialise la donnée émise sur MOSI, et retourne la donnée reçue via MISO.\\
	{\tt byte echangeSPI (const byte inDonneeEmise) ;}

	\item {\bf Écriture d'un registre :} Affecte inValeur au registre {\tt inAdresse}.\\
	{\tt void ecrireRegistreMCP2515 (const byte inAddresse, const byte inValeur) ;}

	\item {\bf Lecture d'un registre :} Cette fonction retourne la valeur du registre {\tt inAdresse.}\\
	{\tt  byte lireRegistreMCP2515 (const byte inAddresse) ;}
\end{itemize}
\bigskip

L'envoie des bits se fait en MSB et il ne faut pas remettre à la ligne {\tt SS} entre la transmission de 2 octets.
Vous pouvez utiliser {\tt delayMicroseconds} (\href{https://www.arduino.cc/reference/en/language/functions/time/delaymicroseconds/}{documentation}) pour créer le délai.

{\bf Comment savoir si l'écriture et la lecture fonctionne correctement ?} On va écrire une valeur dans un registre, et relire ce registre et comparer la valeur lue avec la valeur écrite. Il faut sélectionner un registre qui se comporte comme une mémoire : par exemple, le registre TXB0D0, d'adresse 0x36. {\bf Ecrivez votre fonction de test dans le croquis Arduino et faites une capture d'écran avec l'analyseur logique pour montrer les trames.}


%%%%%%%%%%%%%%%%%%%%%%%%%
  \section{Utiliser le drive Arduino}
%%%%%%%%%%%%%%%%%%%%%%%%%

	Vous allez programmer une liaison SPI matérielle pour communiquer avec le MCP2515. Pour utiliser la librairie Arduino, il faut commencer par configurer les broches d'entrée-sortie convenablement. Pour cela suivez les indications ci-dessous.

{\bf Inclusions :} Deux fichiers d'en-tête sont à include :\\
{\tt \#include <SPI.h>}\\
{\tt \#include <wiring\_private.h>}\\
\bigskip

{\bf Variable globale :} Une variable globale mySPI est à déclarer :\\
{\tt SPIClass mySPI (\&sercom1, 10, 12, 11, SPI\_PAD\_0\_SCK\_3, SERCOM\_RX\_PAD\_2) ;}
\bigskip

{\bf Initialisation :} Écrire une fonction initialiserSPIMateriel dans laquelle :
\begin{enumerate}
	\item le port correspondant à SS est configuré en sortie, et placé au niveau haut (comme pour le SPI logiciel) ; 
	\item appeler {\tt mySPI.begin ()}, cet appel configure la liaison SPI avec des paramètres par défaut ;
	\item appeler {\tt mySPI.beginTransaction(SPISettings (10 * 1000 * 1000, MSBFIRST, SPI\_MODE0))}, 
cet appel fixe la vitesse (10 * 1000 * 1000 = 10 MHz), l'ordre de sérialisation des bits, et le mode ; 
	\item appeler {\tt pinPeripheral (10, PIO\_SERCOM)}, cet appel relie le port n°10 au module SERCOM ; 
	\item appeler {\tt pinPeripheral (11, PIO\_SERCOM)}
	\item appeler {\tt pinPeripheral (12, PIO\_SERCOM)}
\end{enumerate}
\bigskip

{\bf Écriture d'un registre, lecture d'un registre :} Appeler la fonction {\tt mySPI.transfer} à la place de la fonction {\tt echangeSPI} utilisée lors de la liaison SPI logicielle.

 {\bf Ecrivez votre fonction de test dans le croquis Arduino et faites une capture d'écran avec l'analyseur logique pour montrer les trames.}


\end{document}